\section*{Chapter \ref{ch:g6:first-python-program} --- A First Python Program}

\begin{solution}{exo:g6:first-python-program:1}
\code{13}: the multiplication $3 \times 2 = 6$ comes first, then $7 + 6$.
\end{solution}

\begin{solution}{exo:g6:first-python-program:2}
\code{5} and \code{2}: $17 = 3 \times 5 + 2$.
\end{solution}

\begin{solution}{exo:g6:first-python-program:3}
Four lines, in the order of the \omterm{def:g4:programs-and-operating-system:program}{program}: \code{3}, \code{2}, \code{1} (the
result of $3 - 2$, computed because there are no quotes), and \code{Go!}.
\end{solution}

\begin{solution}{exo:g6:first-python-program:4}
For example:

\omcode{code/grade-6/01-first-python-program/exercises/poster.py}{1}{3}{A poster
of three lines.}\label{lst:g6:first-python-program:sol-poster}
\end{solution}

\begin{solution}{exo:g6:first-python-program:5}
The closing bracket: \code{print("Hello")}. Every opening bracket needs its
closing bracket.
\end{solution}

\begin{solution}{exo:g6:first-python-program:6}
The last line says that the name \code{Print} is not defined (a
\code{NameError}) and asks whether \code{print} was meant. Python makes a difference between capital and small letters: the
fix is \code{print("Hi")}, with a small p.
\end{solution}

\begin{solution}{exo:g6:first-python-program:7}
\omcode{code/grade-6/01-first-python-program/exercises/weeks.py}{1}{2}{Weeks and
days in 100 days.}\label{lst:g6:first-python-program:sol-weeks}

It prints \code{14 weeks} and \code{2 days}: $100 = 7 \times 14 + 2$.
\end{solution}

\begin{solution}{exo:g6:first-python-program:8}
\code{2.5} and \code{2}. The sign \code{/} gives the exact quotient as a
decimal number; \code{//} gives only the whole number of times 4 goes into 10.
\end{solution}

\begin{solution}{exo:g6:first-python-program:9}
\omcode{code/grade-6/01-first-python-program/exercises/tickets.py}{1}{1}{The
price of 7 tickets.}\label{lst:g6:first-python-program:sol-tickets}

It prints \code{31.5}: the tickets cost 31.50 euros.
\end{solution}

\begin{solution}{exo:g6:first-python-program:10}
Four lines: \code{one}; then \code{three four} (a \code{print} with two things
prints them on one line, with a space between); then an empty line, printed by
\code{print()}; then \code{five}. Line 3 is a \omterm{def:g6:first-python-program:comment}{comment} and prints nothing.
\end{solution}

\begin{solution}{exo:g6:first-python-program:11}
Line 3 computes $50 - 8 = 42$: the cards \emph{not} given to one player, which
means nothing here. The leftover is the remainder of the division:

\omcode{code/grade-6/01-first-python-program/exercises/cards.py}{1}{3}{The
corrected \omterm{def:g4:programs-and-operating-system:program}{program}.}\label{lst:g6:first-python-program:sol-cards}

It prints 8 and 2: $50 = 6 \times 8 + 2$.
\end{solution}

\begin{solution}{exo:g6:first-python-program:12}
The \omterm{def:g4:programs-and-operating-system:program}{program} \code{exercises/give_back.py} prints, for ten numbers, \code{n // 5},
\texttt{n \% 5} and the whole expression; for example 23 gives 4, 3 and 23,
and 137 gives 27, 2 and 137. The expression always gives back $n$. The reason is
the division with a remainder: \code{n // 5} is the number $q$ of whole fives in
$n$, \texttt{n \% 5} is the remainder $r$, and $n$ is exactly $q$ fives plus
$r$: $n = 5 \times q + r$.
\end{solution}

\begin{solution}{pb:g6:first-python-program:1}
\textbf{1.} $137 = 50 \times 2 + 37$: 2 buses are full and 37 pupils are still
waiting, so a third bus is needed: 3 buses.

\textbf{2.} $3 \times 50 = 150$ seats for 137 pupils: 13 seats stay empty.

\textbf{3.} $137 = 12 \times 11 + 5$: 11 groups of 12, and a last group of 5:
12 groups in all.

\textbf{4.} $137 = 50 \times 2 + 37 = 12 \times 11 + 5$. \code{137 // 50} gives 2
and \texttt{137 \% 50} gives 37; \code{137 // 12} gives 11 and \texttt{137 \% 12}
gives 5.

\textbf{5.} The \omterm{def:g4:programs-and-operating-system:program}{program} prints 137 pupils, 3 buses, 13 empty seats and 12
museum groups, as found by hand.

\omoutput{code/grade-6/01-first-python-program/outputs/trip.txt}{The school trip
for 137 pupils.}\label{out:g6:first-python-program:trip}

\textbf{6.} For 150: 3 buses, 0 empty seats, 13 groups. For 151: 4 buses, 49
empty seats, 13 groups. With 150 pupils the three buses are exactly full; one
more pupil needs a fourth bus, almost empty.

\textbf{7.} \code{pupils // 50 + 1} is right for 137 (3), but gives 4 for 150,
where 3 buses are enough. It is wrong exactly when the number of pupils is a
multiple of 50: 50, 100, 150, 200, \dots, where it orders one bus too many.

\textbf{8.} Adding $12 - 1 = 11$ pushes any number of pupils that is not a
multiple of 12 past the next multiple, so the last, smaller group is counted;
a multiple of 12 is not pushed past the next one. It is the same trick as
$50 - 1 = 49$ for the buses.

\textbf{9.} The number of buses goes from 1 to 2 at 51 pupils, from 2 to 3 at
101 and from 3 to 4 at 151, and stays 4 up to 200:

\omoutput[1-1,50-53,100-103,150-153,201-201]{code/grade-6/01-first-python-program/outputs/buses_table.txt}{Around
the changes.}\label{out:g6:first-python-program:buses-table}

\textbf{10.} If the number of pupils is a multiple of 50, adding 49 does not
reach the next multiple: $100 + 49 = 149$, and \code{149 // 50} is 2, the exact
number of buses. If it is not a multiple of 50, there is a remainder of at least
1, and adding 49 carries the number past the next multiple: $101 + 49 = 150$
gives 3. Adding 49 means: every pupil left over asks for one more bus; and
$150 + 49 = 199$ still gives 3.

\textbf{11.} No value disagrees:

\omoutput{code/grade-6/01-first-python-program/outputs/buses_check.txt}{The rule
and ``divide and round up'' agree from 1 to 200
pupils.}\label{out:g6:first-python-program:buses-check}

137 pupils need 3 buses; a school needs 2 buses from 51 pupils on, 3 from 101
and 4 from 151 (up to 200).
\end{solution}
